\documentclass[11pt]{article}
\newcommand{\DocTitle}{Implementation and Migration Blueprint}
\newcommand{\DocSubtitle}{Incremental modernization with HTML, APIs, state machines and rollback}
\newcommand{\DocID}{DOS-08}
% design-of-sites-style.tex
% Suite preamble for the Web Design Patterns / Design of Sites suite.
%
% RECONSTRUCTION PROVENANCE: the nine uploaded suite documents load a shared
% preamble with `\input{style.tex}`, but no `style.tex` was present in the
% upload batches or anywhere in this repository's history. This file is a
% documented, suite-specific reconstruction inferred from the interface the
% documents actually use: per-document metadata (\DocTitle, \DocSubtitle,
% \DocID defined before this file is input), the front-matter renderer
% \MakeFront, the provenance marker \Provenance, the callout environments
% groundedbox/modernbox/notebox/decisionbox, and the colors
% Ink/Ocean/Mint/Amber/Faint referenced in document TikZ code. It is NOT the
% recovered original style; the original visual design is unknown.
%
% pdfLaTeX only.
\usepackage[T1]{fontenc}
\usepackage[utf8]{inputenc}
\usepackage{lmodern}
\usepackage{microtype}
\usepackage[letterpaper,margin=0.85in,headheight=15pt]{geometry}
\usepackage[table]{xcolor}
\definecolor{Ink}{HTML}{1E2A36}
\definecolor{Ocean}{HTML}{0F4C72}
\definecolor{Mint}{HTML}{0E7E6B}
\definecolor{Amber}{HTML}{B07A1E}
\definecolor{Faint}{HTML}{F2F7F9}
\usepackage{amsmath,amssymb}
\usepackage{array,booktabs,longtable,tabularx,multirow}
\usepackage{enumitem}
\usepackage[most]{tcolorbox}
\usepackage{listings}
\usepackage{tikz}
\usetikzlibrary{arrows.meta,positioning,shapes.geometric,fit,calc}
\usepackage{needspace}
\usepackage{titlesec}
\usepackage{fancyhdr}
\usepackage{lastpage}
\usepackage{xurl}
\usepackage{hyperref}
\hypersetup{colorlinks=true,linkcolor=Ocean,urlcolor=Mint,citecolor=Ocean,pdftitle={\DocTitle},pdfauthor={Web Design Patterns Engineering Suite}}
\setlength{\parindent}{0pt}
\setlength{\parskip}{5pt plus 1pt}
\setlength{\emergencystretch}{2em}
\renewcommand{\arraystretch}{1.18}
\setlist[itemize]{itemsep=2pt,topsep=3pt,leftmargin=1.5em}
\setlist[enumerate]{itemsep=2pt,topsep=3pt,leftmargin=1.7em}
\titleformat{\section}{\Large\bfseries\color{Ocean}}{\thesection}{0.7em}{}[\vspace{-3pt}\color{Mint}\titlerule]
\titleformat{\subsection}{\large\bfseries\color{Ink}}{\thesubsection}{0.6em}{}
\titleformat{\subsubsection}{\normalsize\bfseries\color{Mint}}{\thesubsubsection}{0.55em}{}
\pagestyle{fancy}
\fancyhf{}
\fancyhead[L]{\sffamily\footnotesize\color{Ocean} Design of Sites | \DocID}
\fancyhead[R]{\sffamily\footnotesize\color{Ink} Web Design Patterns Suite}
\fancyfoot[L]{\sffamily\scriptsize Original interpretive guidance; pattern names used as a taxonomy}
\fancyfoot[R]{\sffamily\scriptsize\thepage\ / \pageref{LastPage}}
\renewcommand{\headrulewidth}{0.3pt}
\lstdefinestyle{dositescode}{basicstyle=\ttfamily\footnotesize,breaklines=true,columns=fullflexible,keepspaces=true,showstringspaces=false,frame=single,rulecolor=\color{Mint},backgroundcolor=\color{Faint},aboveskip=8pt,belowskip=8pt}
\lstset{style=dositescode}
% Callouts. groundedbox marks statements verified against the public
% bibliographic record; modernbox marks newly authored modern engineering
% interpretation; notebox is a study/summary callout; decisionbox records
% decisions and trade-offs.
\newtcolorbox{groundedbox}[1][]{enhanced,breakable,colback=Faint,colframe=Mint,boxrule=0.5pt,arc=1.5mm,title={Grounded in the public record},fonttitle=\bfseries\small,#1}
\newtcolorbox{modernbox}[1][]{enhanced,breakable,colback=white,colframe=Ocean,boxrule=0.5pt,arc=1.5mm,title={Modern engineering interpretation},fonttitle=\bfseries\small,#1}
\newtcolorbox{notebox}[1][]{enhanced,breakable,colback=Faint,colframe=Amber,boxrule=0.5pt,arc=1.5mm,#1}
\newtcolorbox{decisionbox}[1][]{enhanced,breakable,colback=white,colframe=Amber,boxrule=0.5pt,arc=1.5mm,title={Decision record},fonttitle=\bfseries\small,#1}
% Front matter renderer: consumes \DocTitle, \DocSubtitle, \DocID.
\newcommand{\MakeFront}{%
  \thispagestyle{empty}%
  {\color{Mint}\rule{\textwidth}{3pt}}\par\vspace{16pt}
  {\sffamily\small\color{Amber}\bfseries \DocID\par}\vspace{10pt}
  {\Huge\bfseries\color{Ocean} \DocTitle\par}\vspace{12pt}
  {\Large\color{Ink} \DocSubtitle\par}\vspace{14pt}
  {\small\color{Ink} Independent educational and engineering companion to \emph{The Design of Sites} (2nd edition). Original explanatory material; the copyrighted pattern text is not reproduced.\par}
  \vspace{10pt}{\color{Mint}\rule{\textwidth}{1pt}}\par\vspace{8pt}
  \tableofcontents\clearpage}
% Provenance marker used at the top of each document body.
\newcommand{\Provenance}{\par\noindent{\footnotesize\sffamily\color{Amber}\bfseries Provenance:}~{\footnotesize\color{Ink}\itshape claims grounded in the public bibliographic record are boxed and cited; everything else is newly authored interpretation.}\par\smallskip}

\begin{document}
\MakeFront
\section{Scenario and objective}

\Provenance
\begin{groundedbox}The historical reference categories used here are B/C/D/F/H/K/L/M [B1]. The hypothetical application and all technical designs/code in this document are newly created teaching examples.\end{groundedbox}
\textbf{Scenario.} A legacy commerce portal uses a server-rendered catalog, several client-side widgets, desktop-only navigation and a monolithic checkout endpoint. The goal is an accessible, secure and measurable experience without an avoidable rewrite.
\textbf{Non-goals.} Framework migration for its own sake, one microservice per pattern, replacing usable server rendering with unnecessary hydration, and collecting additional personal data to increase analytics volume.

\section{As-is assessment checklist}
\begin{itemize}
\item Inventory top user journeys, conversion interruptions, core route graph and customer support issues.
\item Inventory browser entry points, pages, forms, identity/session flows, static asset pipeline and APIs.
\item Assess semantic HTML, keyboard operation, 400 percent zoom, focus and screen reader critical tasks.
\item Enumerate trust boundaries, payment-provider integration, price authority and log-data exposure.
\item Capture baseline field web vitals, backend latency, availability and deployment recovery capability.
\item Identify owners and versions of shared UI libraries and unpublished CMS schemas.
\end{itemize}

\section{Strangler-style delivery milestones}

\begin{longtable}{@{}p{20mm}p{40mm}p{64mm}p{27mm}@{}}\toprule
Phase & Delivery target & Exit evidence & Key patterns\\\midrule\endhead
0 Discover & Research and measurement foundation & Task baselines, threat and route inventory & A/B/E\\
1 Stabilize & Semantic templates and error recovery & Accessibility and E2E smoke tests & D/H/I/K\\
2 Extract & Catalog and search contract & Versioned API, ACL tests, query relevance & B/D/J\\
3 Harden & Cart/order state and identity & Replay, cross-user, payment exceptions & E/F/H\\
4 Optimize & CDN and rendering/profile changes & Device-segment Web Vitals and tracing & L/M\\
5 Evolve & Ownership boundaries if justified & ADR, release independence and rollback & A/F/G\\
\bottomrule\end{longtable}
\textit{Phase numbering is an implementation plan, not the historical chapter order.}

\section{HTML-first accessible form example}

\begin{lstlisting}[language=HTML]
<form action="/checkout/address" method="post">
  <h2>Delivery address</h2>
  <label for="street">Street address</label>
  <input id="street" name="street"
         autocomplete="shipping street-address" required>
  <p id="street-help">Use the address where you can
     receive deliveries.</p>
  <button type="submit">Continue to shipping</button>
</form>
\end{lstlisting}
\textbf{Important:} This shows a minimum semantic baseline, not a complete checkout implementation. Production must associate help using \texttt{aria-describedby}, present validation errors programmatically, use CSRF defenses as applicable and implement server-side validation and authorization.

\section{Progressively enhanced task state}

\begin{enumerate}
\item GET a canonical route to load state as accessible HTML from the server.
\item POST a transition with server-side validation, CSRF protection where applicable and idempotency as required.
\item Redirect after successful mutation to avoid accidental form resubmission where appropriate.
\item Add client enhancements only after the basic journey and recovery path work; announce dynamic changes accessibly.
\item Store only necessary draft state and provide explicit expiration behavior.
\end{enumerate}
\begin{modernbox}A state machine is preferable to a collection of implicit flags. For checkout, use states such as \texttt{Draft}, \texttt{Reviewed}, \texttt{Submitted}, \texttt{Authorized}, \texttt{Confirmed}, \texttt{Failed} and \texttt{Cancelled}. Define allowed commands, retries and side-effect boundaries for each transition.\end{modernbox}

\section{Server-side idempotency and authorization sketch}

\begin{lstlisting}[language={}]
// Illustrative pseudocode: one transaction and durable key scope.
function confirmOrder(actor, cartId, expectedVersion, key) {
  authorize(actor, "checkout:create", cartId);
  return transaction(() => {
    const existing = commands.find(actor.id, "checkout", key);
    if (existing) return existing.persistedResult;
    const cart = carts.lockAndRead(cartId);
    if (cart.version !== expectedVersion) throw Conflict();
    const price = pricing.recalculate(cart.items);
    requireExplicitApproval(actor, cart, price);
    const order = orders.createFromApprovedCart(cart, price);
    commands.save(actor.id, "checkout", key, order.id);
    outbox.enqueue("OrderCreated", order.id);
    return { orderId: order.id };
  });
}
\end{lstlisting}
\textbf{Production caveat:} The example omits implementation-specific payment authorization, storage retention, distributed command semantics and retry compensation. Never treat \texttt{actor.id} from an unverified client request as proof of identity; use authenticated server context. A persistent uniqueness constraint is required for idempotency.

\section{Component ownership and contracts}

\begin{tabularx}{\linewidth}{@{}p{30mm}XX@{}}\toprule
Contract & Producer obligation & Consumer obligation\\\midrule
Design token package & Version semantics, contrast-safe defaults, documentation & Test visual and forced-colors variants\\
Navigation shell & Landmarks, identity state, route fallback, keyboard behavior & Do not override global focus semantics\\
Catalog API & Stable identifiers, availability and cache policy & Display freshness and handle unavailable variants\\
Checkout API & Authoritative pricing, idempotency, traceable errors & Avoid optimistic success before confirmation\\
Telemetry schema & Redaction rules, units, version, retention & No credentials, private query or payment logging\\
\bottomrule\end{tabularx}

\section{Release and rollback choreography}

\begin{enumerate}
\item Introduce route-based traffic segmentation and capture the baseline task completion rate.
\item Launch read-only pages first; compare SEO, performance, navigation and accessible behavior.
\item Use contract tests and dual-run shadow reads where safe. Do not duplicate state-changing payment calls.
\item Canary checkout by a controlled cohort with immediate disablement and state reconciliation procedures.
\item Verify the order/event ledger before increasing traffic and audit cross-user authorization failures.
\item Remove legacy paths only after bookmarks, redirects and operational runbooks are validated.
\end{enumerate}
\begin{decisionbox}Rollback must account for \emph{data and workflow state}, not only code. An old frontend must not silently misinterpret an order created using a newer state model.\end{decisionbox}

\section{Operational runbook example}

\textbf{Incident:} Order confirmation page times out after payment authorization.
\begin{enumerate}
\item Identify correlation by opaque order ID, not customer email or card details.
\item Inspect checkout transition and provider webhook event; do not blindly resubmit a payment request.
\item Query idempotency ledger and outbox for durable commit state.
\item If order exists, return a stable confirmation route; if state is indeterminate, show a recovery message and route to human support.
\item Reconcile asynchronously with provider records under dual controls and audit changes.
\item Review user-facing error language, alerts and synthetic recovery tests.
\end{enumerate}


\clearpage
% design-of-sites-references.tex
% Shared references section for the Web Design Patterns / Design of Sites
% suite.
%
% RECONSTRUCTION PROVENANCE: the nine uploaded suite documents end with
% `\input{references.tex}`, but no `references.tex` was present in the upload
% batches or anywhere in this repository's history. The documents cite only
% keys [B1] and [B2] (bracketed tokens such as "[B4, F1, H10, ...]" are
% pattern-ID template placeholders, not citations). The two entries below are
% reconstructed from the bibliographic facts the documents themselves state:
% publisher record and bibliographic-database record of the second edition.
% This file is NOT the recovered original reference list.
\clearpage
\section*{References}
\begin{description}[leftmargin=3.2em,style=nextline]
  \item[{[B1]}] Publisher bibliographic record and public table of contents:
    D.~K.~van Duyne, J.~A.~Landay, J.~I.~Hong, \emph{The Design of Sites:
    Patterns, Principles, and Processes for Crafting a Customer-Centered Web
    Experience}, 2nd edition, Prentice Hall, December 2006, 1{,}024 pages.
    Confirms the five foundational chapters, the thirteen pattern groups
    (A--M), the 107 pattern labels and the research appendixes used as this
    suite's taxonomy.
  \item[{[B2]}] Bibliographic-database record (O'Reilly learning-platform
    catalog entry) for the same second edition; corroborates edition,
    authors, publication date and chapter/group structure. The copyrighted
    pattern explanations themselves were not retrieved or reproduced for
    this suite.
\end{description}

\end{document}
